% !TeX root = ../eccv2016submission.tex

\begin{figure*}[t]
	\begin{center}
		\vspace{-2ex}
		\centerline{\includegraphics[width=1\columnwidth]{figures/resblock.eps}}
		\vspace{-2ex}
		\caption{A close look at the $\ell^\text{th}$ ResBlock in a ResNet.\label{fig:resblock}}
		\vspace{-6ex}
	\end{center}
\end{figure*}

Learning with \name{} is based on a simple intuition. To reduce the \emph{effective} length of a neural network during training, we randomly skip layers entirely. We achieve this by introducing skip connections in the same fashion as ResNets, however the connection pattern is randomly altered for each mini-batch. For each mini-batch we randomly select sets of layers and remove their corresponding transformation functions, only keeping the identity skip connection.
Throughout, we use the architecture described by He et al.~\cite{he2015deep}. Because the architecture already contains skip connections, it is straightforward to modify, and isolates the benefits of \name{} from that of the ResNet identity connections. Next we describe this network architecture and then explain the \name{} training procedure in detail.

\paragraph{\textbf{ResNet architecture.}}
Following He et al.~\cite{he2015deep}, we construct our network as the functional composition of $L$ \emph{residual blocks} (ResBlocks), each encoding the update rule~\eqref{eq:res_block}.  Fig.~\ref{fig:resblock} shows a schematic illustration of the $\ell^{th}$ ResBlock. In this example, $f_\ell$ consists of a sequence of layers: {\tt Conv-BN-ReLU-Conv-BN}, where {\tt Conv} and {\tt BN} stand for Convolution and Batch Normalization respectively.
This construction scheme is adopted in all our experiments except ImageNet, for which we use the bottleneck block detailed in He et al. \cite{he2015deep}. Typically, there are 64, 32, or 16 filters in the convolutional layers (see Section~\ref{sec_results} for experimental details).

\begin{figure*}[t]
	\begin{center}
		\vspace{-2ex}
		\centerline{\includegraphics[width=1\columnwidth]{figures/resnet_survival.eps}}
		\caption{The linear decay of $p_\ell$ illustrated on a ResNet with \name{} for $p_0\!=\!1$ and $p_L\!=\! 0.5$. Conceptually, we treat the input to the first ResBlock as $H_0$, which is always active.}
    \label{figure.lin_decay}
		\vspace{-6ex}
	\end{center}
\end{figure*}

\paragraph{\textbf{Stochastic depth}} aims to shrink the depth of a network during training, while keeping it unchanged during testing. We can achieve this goal by randomly dropping entire ResBlocks during training and bypassing their transformations through skip connections.
Let $b_\ell\in\{0,1\}$ denote a Bernoulli random variable, which indicates whether the $\ell^{th}$ ResBlock is active ($b_\ell=1$) or inactive ($b_\ell=0$). Further, let us denote the ``survival'' probability of ResBlock $\ell$ as $p_\ell=\Pr(b_\ell=1)$.

With this definition we can bypass the $\ell^{th}$ ResBlock by multiplying its function $f_\ell$ with $b_\ell$ and we extend the update rule from~\eqref{eq:res_block} to
\begin{equation}
\label{res_block_drop}
 H_\ell = {\tt ReLU} (b_\ell f_\ell(H_{\ell-1}) +  \text{id}(H_{\ell-1})).
\end{equation}
If $b_\ell\!=\!1$, eq.~\eqref{res_block_drop} reduces to the original ResNet update \eqref{eq:res_block} and this ResBlock remains unchanged.
If $b_\ell\!=\!0$, the ResBlock reduces to the identity function,
\begin{equation}
 H_\ell = \text{id}(H_{\ell-1}).
\end{equation}
This reduction follows from the fact that the input $H_{\ell-1}$ is always non-negative, at least for the architectures we use. For $\ell\geq 2$, it is the output of the previous ResBlock, which is non-negative because of the final ${\tt ReLU}$ transition function (see Fig.~\ref{fig:resblock}). For $\ell\!=\!1$, its input is the output of a {\tt Conv-BN-ReLU} sequence that begins the architecture before the first ResBlock. For non-negative inputs the ${\tt ReLU}$ transition function acts as an identity.

\paragraph{\textbf{The survival probabilities}} $p_\ell$ are new hyper-parameters of our training procedure. Intuitively, they should take on similar values for neighboring ResBlocks. One option is to set $p_\ell=p_L$ uniformly for all $\ell$ to obtain a single hyper-parameter $p_L$. Another possibility is to set them according to a smooth function of $\ell$. We propose a simple linear decay rule from $p_0=1$ for the input, to $p_L$ for the last ResBlock:
\begin{equation}
\label{eq:surv_rate}
 p_\ell = 1-\frac{\ell}{L}(1-p_L).
\end{equation}
%where $L$ is the total number of ResBlocks.
See Fig.~\ref{figure.lin_decay} for a schematic illustration.
The linearly decaying survival probability originates from our intuition that the earlier layers extract low-level features that will be used by later layers and should therefore be more reliably present.
In Section~\ref{sec_results} we perform a more detailed empirical comparison between the uniform and decaying assignments for $p_\ell$.  We conclude that the linear decay rule \eqref{eq:surv_rate} is preferred and, as training with \name{} is surprisingly stable with respect to $p_L$, we set $p_L=0.5$ throughout (see Fig.~\ref{figure.deathRate}).

\paragraph{\textbf{Expected network depth}.}
During the forward-backward pass the transformation $f_\ell$ is bypassed with probability $(\!1-p_\ell\!)$, leading to a network with reduced depth.
With \name{}, the number of effective ResBlocks during training, denoted as $\tilde{L}$, becomes a random variable. Its expectation is given by:
$
E(\tilde{L}) =  \sum\nolimits_{\ell=1}^L p_\ell.
$

Under the linear decay rule with $p_L=0.5$, the expected number of ResBlocks during training reduces to $E(\tilde{L})=(3L-1)/4$, or $E(\tilde{L})\approx 3L/4$ when $L$ is large. For the 110-layer network with $L=54$ commonly used in our experiments, we have $E(\tilde{L})\approx 40$. In other words, with \name{}, we train ResNets with an average number of 40 ResBlocks, but recover a ResNet with 54 blocks at test time.
This reduction in depth significantly alleviates the vanishing gradients and the information loss problem in deep ResNets. Note that because the connectivity is random, there will be updates with significantly shorter networks and more direct paths to individual layers. We provide an empirical demonstration of this effect in Section \ref{sec_analysis}.

\paragraph{\textbf{Training time savings}.}
When a ResBlock is bypassed for a specific iteration, there is no need to perform forward-backward computation or gradient updates. As the forward-backward computation dominates the training time, \name{} significantly speeds up the training process. Following the calculations above, approximately $25\%$ of training time could be saved under the linear decay rule with $p_L=0.5$. The timings in practice using our implementation are consistent with this analysis (see the last paragraph of Section~\ref{sec_results}). More computational savings can be obtained by switching to a uniform probability for $p_\ell$ or lowering $p_L$ accordingly. In fact, Fig.~\ref{figure.deathRate} shows that with $p_L=0.2$, the ResNet with \name{} obtains the same test error as its constant depth counterpart on CIFAR-10 but gives a 40\% speedup.

\paragraph{\textbf{Implicit model ensemble}.}
In addition to the predicted speedups, we also observe significantly lower testing errors in our experiments, in comparison with ResNets of constant depth.
One explanation for our performance improvements is that training with \name{} can be viewed as training an ensemble of ResNets \emph{implicitly}.
Each of the $L$ layers is either active or inactive, resulting in $2^L$ possible network combinations. For each training mini-batch one of the $2^L$ networks (with shared weights) is sampled and updated. During testing all networks are averaged using the approach in the next paragraph.

\paragraph{\textbf{Stochastic depth during testing}} requires small modifications to the network. We keep all functions $f_\ell$ active throughout testing in order to utilize the full-length network with all its model capacity. However, during training, functions $f_\ell$ are only active for a fraction $p_\ell$ of all updates, and the corresponding weights of the next layer are calibrated for this survival probability. We therefore need to re-calibrate the outputs of any given function $f_\ell$ by the expected number of times it participates in training, $p_\ell$. The forward propagation update rule becomes:
\begin{equation}
\label{res_block_drop_test}
 H_\ell^\text{Test} = {\tt ReLU} (p_\ell f_\ell(H_{\ell-1}^\text{Test}; W_\ell) +  H_{\ell-1}^\text{Test}).
\end{equation}
From the model ensemble perspective, the update rule~\eqref{res_block_drop_test} can be interpreted as combining all possible networks into a single test architecture, in which each layer is weighted by its survival probability.
